% Tradução pt-BR da passagem espanhola congelada; a fonte permanece somente leitura.
% Tradução da matéria preliminar das páginas 1--9 do escaneamento original.
% O rótulo da ramificação refere-se às leituras da fonte; o português permanece idiomático.
\hypersetup{pageanchor=false}

% SGA-SEG BEGIN T2-FM-000001 kind=scan_page scan_page=1
\SGAFrontCover
  {Lecture Notes in Mathematics}
  {Uma coleção de relatórios e seminários informais}
  {Editada por A. Dold, Heidelberg, e B. Eckmann, Zurique}
  {152}
  {Grupos de tipo multiplicativo\\ e estrutura dos\\ esquemas de grupos\\ em geral}
  {Springer-Verlag}
  {Berlin · Heidelberg · New York 1970}
% O emblema gráfico da Springer não é reproduzido nesta transcrição tipográfica.
% SGA-SEG END T2-FM-000001

% SGA-SEG BEGIN T2-FM-000002 kind=scan_page scan_page=2
\SGACopyrightPage{%
Esta obra é protegida por direitos de autor. Reservam-se todos os direitos, tanto sobre a totalidade como sobre qualquer parte do material, em particular sobre a tradução, a reimpressão, a reutilização de ilustrações, a radiodifusão, a reprodução por fotocopiadora ou meios similares e o armazenamento em bancos de dados.

De acordo com o § 54 da Lei alemã de direitos autorais, quando forem feitas cópias para uso diferente do privado, deverá ser paga uma taxa à editora; o valor dessa taxa será determinado de comum acordo com a editora.

© Springer-Verlag Berlin · Heidelberg 1970. Número de ficha do catálogo da Biblioteca do Congresso: 75-134652. Impresso na Alemanha. N.º de título 3309.

Impressão offset: Julius Beltz, Weinheim/Bergstr.

\medskip
\noindent\textit{Nota editorial: este aviso é conservado unicamente como registro histórico de proveniência e não restringe esta edição matemática independente nem sua publicação.}}
% SGA-SEG END T2-FM-000002

% SGA-SEG BEGIN T2-FM-000003 kind=scan_page scan_page=3
\SGAInnerTitle
  {SEMINÁRIO DE GEOMETRIA ALGÉBRICA DO BOIS MARIE}
  {1962/64}
  {ESQUEMAS DE GRUPOS\\(SGA 3)}
  {Um seminário dirigido por\\[0.7em] M. DEMAZURE e A. GROTHENDIECK\\[1.5em] com a colaboração de\\[0.7em] M. ARTIN, J. E. BERTIN, P. GABRIEL, M. RAYNAUD, J. P. SERRE}
  {TOMO II\\[0.8em]\normalsize (GRUPOS DE TIPO MULTIPLICATIVO E ESTRUTURA DOS ESQUEMAS DE GRUPOS EM GERAL)}
% SGA-SEG END T2-FM-000003

\hypersetup{pageanchor=true}
\pagenumbering{Roman}
\setcounter{page}{4}

% SGA-SEG BEGIN T2-FM-000004 kind=scan_page scan_page=4
\begin{SGAContentsPage}{ÍNDICE GERAL}
\thispagestyle{empty}
\SGAContentsExpose{EXPOSIÇÃO VIII}{GRUPOS DIAGONALIZÁVEIS, por A. GROTHENDIECK.}{1}
\SGAContentsEntry{1.}{Bidualidade}{1}
\SGAContentsEntry{2.}{Propriedades esquemáticas dos grupos diagonalizáveis}{6}
\SGAContentsEntry{3.}{Propriedades de exatidão do funtor $D_S$}{7}
\SGAContentsEntry{4.}{Torsores sob um grupo diagonalizável}{11}
\SGAContentsEntry{5.}{Quociente de um esquema afim por um grupo diagonalizável que atua livremente}{15}
\SGAContentsEntry{6.}{Morfismos essencialmente livres e representabilidade de certos funtores da forma $\prod_{Y/S} Z/Y$}{20}
\SGAContentsEntry{7.}{Apêndice: sobre os monomorfismos de pré-esquemas de grupos}{25}

\SGAContentsExpose{EXPOSIÇÃO IX}{GRUPOS DE TIPO MULTIPLICATIVO: HOMOMORFISMOS EM UM ESQUEMA DE GRUPOS, por A. GROTHENDIECK.}{37}
\SGAContentsEntry{1.}{Definições}{37}
\SGAContentsEntry{2.}{Extensão de certas propriedades dos grupos diagonalizáveis aos grupos de tipo multiplicativo}{40}
\SGAContentsEntry{3.}{Propriedades infinitesimais: teorema de levantamento e conjugação}{46}
\SGAContentsEntry{4.}{O teorema da densidade}{50}
\SGAContentsEntry{5.}{Homomorfismos centrais dos grupos de tipo multiplicativo}{58}
\SGAContentsEntry{6.}{Monomorfismos de grupos de tipo multiplicativo e fatoração canônica de um homomorfismo de um desses grupos}{63}
\SGAContentsEntry{7.}{Algebraicidade dos homomorfismos formais em um grupo afim}{68}
\SGAContentsEntry{8.}{Subgrupos, grupos quocientes e extensões de grupos de tipo multiplicativo sobre um corpo}{74}
\end{SGAContentsPage}
% SGA-SEG END T2-FM-000004

% SGA-SEG BEGIN T2-FM-000005 kind=scan_page scan_page=5 printed_folio=V
\begin{SGAContentsPage}{}
\SGAContentsExpose{EXPOSIÇÃO X}{CARACTERIZAÇÃO E CLASSIFICAÇÃO DOS GRUPOS DE TIPO MULTIPLICATIVO, por A. GROTHENDIECK.}{77}
\SGAContentsEntry{1.}{Classificação dos grupos isotriviais. Caso de um corpo-base}{77}
\SGAContentsEntry{2.}{Variações infinitesimais da estrutura}{81}
\SGAContentsEntry{3.}{Variações de estruturas finitas: anel-base completo}{86}
\SGAContentsEntry{4.}{Caso de uma base arbitrária. Teorema de quase-isotrivialidade}{92}
\SGAContentsEntry{5.}{Esquema de homomorfismos de um grupo de tipo multiplicativo em outro. Grupos constantes torcidos e grupos de tipo multiplicativo}{98}
\SGAContentsEntry{6.}{Recobrimentos principais galoisianos infinitos e grupo fundamental ampliado}{106}
\SGAContentsEntry{7.}{Classificação dos pré-esquemas constantes torcidos e dos grupos de tipo multiplicativo de tipo finito em termos do grupo fundamental ampliado}{112}
\SGAContentsEntry{8.}{Apêndice: eliminação de certas hipóteses afins}{116}

\SGAContentsExpose{EXPOSIÇÃO XI}{CRITÉRIOS DE REPRESENTABILIDADE. APLICAÇÕES AOS SUBGRUPOS DE TIPO MULTIPLICATIVO DOS ESQUEMAS DE GRUPOS AFINS, por A. GROTHENDIECK.}{127}
\SGAContentsEntry{0.}{Introdução}{127}
\SGAContentsEntry{1.}{Recordação sobre os morfismos lisos, étale e não ramificados}{128}
\SGAContentsEntry{2.}{Exemplos de funtores formalmente lisos provenientes da teoria dos grupos de tipo multiplicativo}{137}
\SGAContentsEntry{3.}{Resultados auxiliares de representabilidade}{141}
\SGAContentsEntry{4.}{O esquema dos subgrupos de tipo multiplicativo de um grupo liso afim}{157}
\SGAContentsEntry{5.}{Primeiros corolários do teorema de representabilidade}{164}
\SGAContentsEntry{6.}{Sobre uma propriedade de rigidez para os homomorfismos de certos esquemas de grupos e a representabilidade de certos transportadores}{171}
\end{SGAContentsPage}
% SGA-SEG END T2-FM-000005

% SGA-SEG BEGIN T2-FM-000006 kind=scan_page scan_page=6 printed_folio=VI
\begin{SGAContentsPage}{}
\SGAContentsExpose{EXPOSIÇÃO XII}{TOROS MAXIMAIS, GRUPO DE WEYL, SUBGRUPO DE CARTAN, CENTRO REDUTIVO DOS ESQUEMAS DE GRUPOS LISOS E AFINS, por A. GROTHENDIECK.}{180}
\SGAContentsEntry{1.}{Toros maximais}{181}
\SGAContentsEntry{2.}{O grupo de Weyl}{191}
\SGAContentsEntry{3.}{Subgrupos de Cartan}{196}
\SGAContentsEntry{4.}{O centro redutivo}{198}
\SGAContentsEntry{5.}{Aplicação ao esquema dos subgrupos de tipo multiplicativo}{210}
\SGAContentsEntry{6.}{Toros maximais e subgrupos de Cartan dos grupos algébricos não necessariamente afins (corpo-base algebricamente fechado)}{216}
\SGAContentsEntry{7.}{Aplicação aos pré-esquemas de grupos lisos não necessariamente afins}{224}
\SGAContentsEntry{8.}{Elementos semissimples, união e interseção dos toros maximais nos esquemas de grupos não necessariamente afins}{239}

\SGAContentsExpose{EXPOSIÇÃO XIII}{ELEMENTOS REGULARES DOS GRUPOS ALGÉBRICOS E DAS ÁLGEBRAS DE LIE, por A. GROTHENDIECK.}{249}
\SGAContentsEntry{1.}{Um lema auxiliar sobre as variedades com operadores}{249}
\SGAContentsEntry{2.}{Teorema da densidade e teoria dos pontos regulares de $G$}{253}
\SGAContentsEntry{3.}{Caso de um pré-esquema-base arbitrário}{269}
\SGAContentsEntry{4.}{Álgebras de Lie sobre um corpo: posto, elementos regulares, subálgebras de Cartan}{276}
\SGAContentsEntry{5.}{Caso da álgebra de Lie de um grupo algébrico liso: teorema da densidade}{283}
\SGAContentsEntry{6.}{Subálgebras de Cartan e subgrupo de tipo (C), relativos a um grupo algébrico liso}{291}

\SGAContentsExpose{EXPOSIÇÃO XIV}{ELEMENTOS REGULARES: CONTINUAÇÃO. APLICAÇÕES AOS GRUPOS ALGÉBRICOS, por A. GROTHENDIECK; APÊNDICE por J. P. SERRE.}{296}
\SGAContentsEntry{1.}{Construção de subgrupos de Cartan e de toros maximais para um grupo algébrico liso}{296}
\SGAContentsEntry{2.}{Álgebras de Lie sobre um pré-esquema arbitrário: seções regulares e subálgebras de Cartan}{299}
\end{SGAContentsPage}
% SGA-SEG END T2-FM-000006

% SGA-SEG BEGIN T2-FM-000007 kind=scan_page scan_page=7 printed_folio=VII
\begin{SGAContentsPage}{}
% Continuação da EXPOSIÇÃO XIV.
\SGAContentsEntry{3.}{Subgrupo de tipo (C) dos pré-esquemas de grupos sobre um pré-esquema arbitrário}{312}
\SGAContentsEntry{4.}{Uma digressão sobre os subgrupos de Borel}{323}
\SGAContentsEntry{5.}{Relações entre subgrupos de Cartan e subálgebras de Cartan}{330}
\SGAContentsEntry{6.}{Aplicações à estrutura dos grupos algébricos}{334}
\SGAContentsEntry{7.}{Apêndice: existência de elementos regulares sobre corpos finitos}{342}

\SGAContentsExpose{EXPOSIÇÃO XV}{COMPLEMENTOS SOBRE OS SUBTOROS DE UM PRÉ-ESQUEMA DE GRUPOS. APLICAÇÃO AOS GRUPOS LISOS, por M. RAYNAUD.}{349}
\SGAContentsEntry{0.}{Introdução}{349}
\SGAContentsEntry{1.}{Levantamento dos subgrupos finitos}{350}
\SGAContentsEntry{2.}{Levantamento infinitesimal dos subtoros}{357}
\SGAContentsSubentry{1.}{Enunciado do teorema}{357}
\SGAContentsSubentry{2.}{Demonstração de 2.1}{362}
\SGAContentsEntry{3.}{Caracterização de um subtoro por seu conjunto subjacente}{374}
\SGAContentsSubentry{1.}{Enunciado do teorema}{374}
\SGAContentsSubentry{2.}{Demonstração das afirmações «fáceis» contidas em 3.1}{377}
\SGAContentsSubentry{3.}{Continuação da demonstração de 3.1}{380}
\SGAContentsEntry{4.}{Caracterização de um subtoro $T$ por meio dos subgrupos ${}_nT$}{398}
\SGAContentsSubentry{1.}{Enunciado do teorema principal}{398}
\SGAContentsSubentry{2.}{Aplicação}{400}
% SGA-SRC-000050: o francês diplomático funde sous-groupeslisses; o português permanece idiomático.
\SGAContentsEntry{5.}{Representabilidade do funtor: subgrupos lisos iguais ao seu normalizador conexo}{409}
\SGAContentsEntry{6.}{Funtor dos subgrupos de Cartan e funtor dos subgrupos parabólicos}{422}
\SGAContentsEntry{7.}{Subgrupos de Cartan de um grupo liso}{445}
\SGAContentsEntry{8.}{Critério de representabilidade do funtor dos subtoros de um grupo liso}{459}
\end{SGAContentsPage}
% SGA-SEG END T2-FM-000007

% SGA-SEG BEGIN T2-FM-000008 kind=scan_page scan_page=8 printed_folio=VIII
\begin{SGAContentsPage}{}
\SGAContentsExpose{EXPOSIÇÃO XVI}{GRUPOS DE POSTO UNIPOTENTE NULO, por M. RAYNAUD.}{484}
\SGAContentsEntry{1.}{Um critério de imersão}{484}
\SGAContentsEntry{2.}{Um teorema de representabilidade dos quocientes}{503}
\SGAContentsEntry{3.}{Grupos com centro plano}{510}
\SGAContentsEntry{4.}{Grupos com fibras afins, de posto unipotente nulo}{520}
\SGAContentsEntry{5.}{Aplicação aos grupos redutivos e semissimples}{524}
\SGAContentsEntry{6.}{Aplicações: extensão de certas propriedades de rigidez dos toros aos grupos de posto unipotente nulo}{527}

\SGAContentsExpose{EXPOSIÇÃO XVII}{GRUPOS ALGÉBRICOS UNIPOTENTES. EXTENSÕES ENTRE GRUPOS UNIPOTENTES E GRUPOS DE TIPO MULTIPLICATIVO, por M. RAYNAUD.}{532}
\SGAContentsEntry{0.}{Algumas notações}{532}
\SGAContentsEntry{1.}{Definição dos grupos algébricos unipotentes}{534}
\SGAContentsEntry{2.}{Primeiras propriedades dos grupos unipotentes}{538}
\SGAContentsEntry{3.}{Grupos unipotentes que atuam sobre um espaço vetorial}{543}
\SGAContentsEntry{4.}{Uma caracterização dos grupos unipotentes}{557}
\SGAContentsSubentry{4.1.}{Grupos algébricos lisos, conexos e afins}{557}
\SGAContentsSubentry{4.2.}{Grupos radiciais}{562}
\SGAContentsSubentry{4.3.}{Grupos afins conexos em característica $p>0$}{565}
\SGAContentsSubentry{4.4.}{Grupos étales}{569}
\SGAContentsSubentry{4.5.}{Variedades abelianas}{570}
\SGAContentsSubentry{4.6.}{Caso geral}{570}
\SGAContentsEntry{5.}{Extensão de um grupo de tipo multiplicativo por um grupo unipotente}{573}
\SGAContentsSubentry{5.1.}{Enunciado do teorema}{573}
\SGAContentsSubentry{5.2.}{Demonstração de 5.1.1 i) e ii) no caso em que $U$ é liso e $H$ é étale}{575}
\SGAContentsSubentry{5.3.}{Estudo do caso em que $H$ é liso}{580}
\SGAContentsSubentry{5.4.}{Estudo do caso em que $U$ é radicial}{582}
\SGAContentsSubentry{5.5.}{Demonstração de 5.1.1 i)}{583}
% O original imprime 5.6 sem título.
\SGAContentsSubentry{5.6.}{}{586}
\SGAContentsSubentry{5.7.}{Demonstração de 5.1.1 ii) b)}{588}
\SGAContentsSubentry{5.8.}{Final da demonstração de 5.1.1 i)}{593}
\SGAContentsSubentry{5.9.}{Contraexemplos}{597}
\end{SGAContentsPage}
% SGA-SEG END T2-FM-000008

% SGA-SEG BEGIN T2-FM-000009 kind=scan_page scan_page=9 printed_folio=IX
\begin{SGAContentsPage}{}
% Continuação da EXPOSIÇÃO XVII.
\SGAContentsEntry{6.}{Extensão de um grupo unipotente por um grupo de tipo multiplicativo}{602}
\SGAContentsSubentry{6.1.}{Enunciado do teorema}{602}
\SGAContentsSubentry{6.2.}{Demonstração de 6.1.1 A)}{602}
\SGAContentsSubentry{6.3.}{Demonstração de 6.1.1 B) e C)}{607}
\SGAContentsSubentry{6.4.}{Exemplos de extensões de um grupo unipotente $U$ por um grupo de tipo multiplicativo $H$ que não são triviais}{608}
\SGAContentsEntry{7.}{Grupos algébricos afins nilpotentes}{611}
\SGAContentsSubentry{7.1.}{Extensões de grupos de tipo multiplicativo}{611}
\SGAContentsSubentry{7.2.}{Estrutura dos grupos algébricos afins comutativos}{613}
\SGAContentsSubentry{7.3.}{Estrutura dos grupos algébricos afins nilpotentes}{615}

\SGAContentsAppendix{APÊNDICE I —}{Cohomologia de Hochschild e extensões de grupos algébricos}{618}
\SGAContentsEntry{1.}{Definição dos grupos de cohomologia}{618}
\SGAContentsEntry{2.}{O grupo $\operatorname{Ext}_{\mathrm{alg}}(G,A)$}{619}
\SGAContentsEntry{3.}{Comparação de $H^2(G,A)$ e $\operatorname{Ext}_{\mathrm{alg}}(G,A)$}{621}

\SGAContentsAppendix{APÊNDICE II —}{Revisões e complementos sobre os grupos radiciais}{623}
\SGAContentsEntry{1.}{O morfismo de Frobenius}{623}
\SGAContentsEntry{2.}{Grupos e $p$-álgebras de Lie}{624}
\SGAContentsEntry{3.}{Grupos radiciais e grupos lisos}{625}

\SGAContentsAppendix{APÊNDICE III —}{Observações e complementos às exposições XV, XVI e XVII}{627}

\SGAContentsExpose{EXPOSIÇÃO XVIII}{TEOREMA DE WEIL SOBRE A CONSTRUÇÃO DE UM GRUPO A PARTIR DE UMA LEI RACIONAL, por M. ARTIN.}{632}
\SGAContentsEntry{0.}{Introdução}{632}
\SGAContentsEntry{1.}{«Revisões» sobre as aplicações racionais}{633}
\SGAContentsEntry{2.}{Determinação local de um morfismo de grupos}{635}
\SGAContentsEntry{3.}{Construção de um grupo a partir de uma lei racional}{639}

\SGAContentsIndex{ÍNDICE DE NOTAÇÕES}{654}
\end{SGAContentsPage}
% SGA-SEG END T2-FM-000009
